% sections/supp_methods.tex -- S1 Appendix, part A: supplementary methods (detail moved from the main text so that the
% main text stays readable; nothing needed for reproduction is dropped). Fixed settings checked against docs/protocol.md,
% the configs and the code; measured values are macros.
\section*{A. Supplementary methods}

\subsection*{A1. Data audit}
\emph{Duplicates.} Byte-identical image files were found by SHA-256 hashes of all \NImages{} files. Near-duplicate eyes were
found by representing each mask as ``any vessel'' in 16 $\times$ 16-pixel blocks (a 128 $\times$ 128 grid), dilating it by
one block and computing the intersection over union (IoU) of every pair of masks; pairs with IoU $\geq$ 0.8 were treated
as candidate duplicates.
The pair of a training and a test image with IoU \SameEyeIoU{} shows the same eye.

\emph{Burned-in markings.} Pixels with $\max(R, G, B) > 40$ lying more than 20 pixels outside a preliminary FOV estimate
were merged into regions (closing with a 41-pixel element); a region of at least 20 such pixels was counted as a marking.
For the grouping of possibly same-patient images, each marking was reduced to a binary pattern ($\max(R, G, B) > 128$
outside the FOV, top-left 128 $\times$ 128 pixels, on a 64 $\times$ 64 grid) and patterns were compared by IoU, maximized
over shifts of up to two cells. Patterns were kept in memory only; the text was never read, transcribed or stored. Pairs
with IoU $\geq$ 0.8 were grouped if their numbers were consecutive within the same official split, or if they formed an
isolated pair (a group of exactly two images with that pattern), the case most likely to be one patient. Identical
patterns shared by larger groups were not used, because fragments of common characters also match across patients.

\emph{Split.} Groups were visited in one seeded random order (seed \SplitSeed{}); a group was assigned to the validation set
only if it fitted the remaining validation places of every category it contains, so that each category has exactly
\NValPerCategory{} validation images. That no group crosses the split was checked when the split was written and again
from the saved split files.

\subsection*{A2. Field-of-view mask}
For each image: (1) $g = \max(R, G, B)$, smoothed with a Gaussian of $\sigma = 1$ pixel; (2) threshold at the 99th
percentile of $g$ in the three image corners that carry no text, plus 1.5; (3) morphological opening with 15 pixels,
largest connected component, holes filled (the eye region); (4) boundary points of the eye region, excluding points on
the image frame; (5) least-squares circle fit, refitted after removing points with $\lvert$residual$\rvert$ greater than
max(3 pixels, three robust standard deviations) until stable; (6) FOV = fitted disk clipped to the image, united with the
eye region. The disk includes dim rims that fall below the threshold, and the union keeps a bright notch at the rim that
contains annotated vessels; burned-in text is not connected to the eye and stays outside. \NFOVTouchFrame{} fitted circles
extend beyond the image frame. The FOV was checked visually on the five darkest images and on images where the eye
touches the frame.

\subsection*{A3. Training-mask resizing}
Thirty training masks (stratified: 8 AMD, 7 DR, 8 glaucoma, 7 normal; random seed 0) were reduced to 512 and 1024 by
nearest-neighbour sampling, block averaging with a threshold of 0.5 or 0.25, max-pooling, and block averaging without a
threshold (soft). Each result was enlarged to 2048 by bilinear interpolation with half-pixel centres, thresholded at 0.5
and scored inside the FOV (\sitab{tab:si_resize}). This measures what each training target retains through the evaluation
path, not how a model trained on it behaves.

\subsection*{A4. Effect of the loss on soft targets}
Free per-pixel logits, without a network, were fitted to the soft 512 masks of four training images (one per category)
for 4,000 steps (Adam, learning rate 0.1 with cosine decay, one image per step), once with binary cross-entropy alone and
once with binary cross-entropy plus soft Dice. With binary cross-entropy alone the fitted probabilities equal the targets.
With the soft Dice term, a target of 0.25 (a one-pixel vessel at 512) is driven to about \LossOptQuarter{} and a target of
0.5 to about \LossOptHalf{}: the Dice term pushes values below roughly half of the image's Dice towards zero and values
above it towards one. This describes the optimum of the loss, not a trained model.

\subsection*{A5. SAM implementation}
LoRA (rank 4) was added to the query and value parts of the fused query--key--value projection in each of the 12 encoder
blocks; one LoRA factor was initialized with Kaiming uniform values and the other with zeros, as in
SAMed~\cite{zhang2023samed}. The LoRA update is added to the frozen projection without a scaling factor (equivalent to
$\alpha = r$) and without dropout. Without prompts, the prompt encoder supplies its empty sparse embedding and its learned
no-mask dense embedding; the mask decoder is called once per sample and returns a single mask (no multi-mask output). Before
training, three optimizer steps on one fixed batch were checked: with fp16 gradient scaling, three steps were applied in
seven iterations; on the first applied step only the zero-initialized LoRA factors received gradients, as expected, and from
the second step on all LoRA factors did; all LoRA parameters and the
96 decoder tensors in use changed; the 24 decoder tensors not in use (the extra mask-token heads and the IoU head, which no
loss reaches) did not change; and every frozen tensor of the image and prompt encoders stayed bit-identical. A short probe of
speed and memory selected a micro-batch of two, the fastest per sample (about 15.7 GB peak allocated memory on the L40S),
accumulated twice to the effective batch of four.

\emph{Tiles.} The tiled model takes the four non-overlapping 1024 $\times$ 1024 crops of the full-resolution enhanced image,
each normalized as a separate SAM input; augmentation is applied per tile. Validation Dice for checkpoint selection is
computed per whole image after stitching its four 256 $\times$ 256 logit grids into one 512 $\times$ 512 grid.

\emph{Overlapping tiles.} Nine 1024 $\times$ 1024 crops at a stride of 512 pixels (the four main crops are among them) are
passed through the same model. Each crop's 256 $\times$ 256 logits are placed on the 512 $\times$ 512 grid at its position
and averaged with the separable triangular weight $w(i) = 1 - \lvert (i + 0.5)/128 - 1 \rvert$, $i = 0, \dots, 255$,
applied along both axes. With 50\% overlap the weights sum to one inside the image, so every interior pixel receives a
convex combination of the crops that cover it and no crop border produces an edge. The blended logits pass through the
sigmoid and are scored like every other model, with the threshold chosen on validation.

\emph{Seam check (validation, seed 0).} The discontinuity was measured on each model's own output grid as the mean absolute
probability difference across the seam lines ($x, y = 1024$ at full resolution), divided by the median difference between
other adjacent rows and columns (about 1 means no seam). Accuracy near the seams was compared with a difference-in-differences:
the difference between the tiled model and a seam-free control (U-Net, averaged over its seeds, or whole-image SAM) in bands
of $\pm$16, 32 and 64 pixels around the seams, minus the same difference in bands around the tile midlines ($x, y = 512$
and 1536), because all models are somewhat less accurate along the central lines of the image. Confidence intervals came
from 10,000 bootstrap resamples of the \NVal{} validation images (\sitab{tab:si_seam}).

\subsection*{A6. Oracle prompts}
For each image, positive points were drawn uniformly from the ground-truth vessel skeleton and negative points uniformly
from FOV background at least 10 pixels (at 2048) from any vessel, with one random permutation per image so that the sets
for $K = 5$, 10 and 20 are nested; the random generator was seeded by a fixed seed combined with the image identifier. All positive and negative
points fell on vessel and background respectively. SAM's predictor was run on the enhanced 1024 image with a single
output mask; logits at 1024 were converted to probabilities and scored like the other models.

\subsection*{A7. Cost measurement}
Inference time was measured for the seed-0 models on one L40S with one image per call over the \NTestImages{} test images,
after 10 warm-up calls, with GPU synchronization before and after each call. The timed span runs from the input tensor on
the GPU to the probability map at the model's output grid, including tiling, stitching or blending and the sigmoid; disk
access and CPU preprocessing are excluded, and the outputs were discarded. The trained models ran with 16-bit mixed
precision, as in training and prediction; the oracle used SAM's default single-precision predictor with $K = 5$. The tiled
model processes its four tiles in one call, whereas the overlapping-tile variant processes its nine tiles in chunks of
three and blends the logits on the CPU, which lowers its peak memory and adds transfer time. Floating-point operations were counted for one
image with PyTorch's \texttt{FlopCounterMode} (one multiply-accumulate counts as two operations). Peak GPU memory is the
maximum allocated memory during the timed loop, including the weights. Training time is the sum of epoch times of the
seed-0 run.

\subsection*{A8. Reproducibility}
Python, NumPy and PyTorch (CPU and GPU) were seeded from the run seed; the data loader's generator was reseeded at every
epoch from the seed and the epoch number, so the image order and the augmentation of an epoch are the same whether or not
a run was interrupted and resumed. cuDNN was set to deterministic mode with benchmarking disabled; some GPU kernels remain
non-deterministic, so repeated runs with the same seed are equal in distribution rather than bit-identical. Checkpoints
store the model, optimizer and gradient-scaler states, the epoch, the best validation score, the early-stopping counter
and all random-number states, and a resumed run continues exactly where it stopped (tested).

\subsection*{A9. Post hoc analyses}
These analyses were added during revision, after the main test evaluation. None predicted any image again or chose
anything on the test set; they re-scored the saved probability maps or re-analysed the saved per-image results.

\textit{Width-stratified centreline precision.} The binary prediction was skeletonized, and each predicted skeleton pixel
inside the FOV was assigned to a width class by the prediction's own distance transform, with the same cut-offs as the
recall (thin $< 3$, medium 3--7, thick $> 7$ pixels at 2048). Precision of a class is the share of its predicted
skeleton pixels that lie within 1 pixel ($3 \times 3$ dilation) of a ground-truth vessel pixel, the mirror image of the
recall. A class without predicted skeleton pixels in an image was left out of that image's mean.

\textit{Spur-pruned thin class.} A medial axis develops short terminal branches (spurs) that run from the centre line of
a wider vessel to small bumps in its boundary; their tips have small radii and fall into the thin class. Terminal
branches, with an end point at one end and a junction at the other, whose pixel count was at most 1.5 times the largest
radius along the branch plus 2 pixels were removed in one pass before width classing, for the ground-truth skeleton
(recall) and the predicted skeleton (precision). This removed \PHPrunedThinShare\% of the ground-truth thin-vessel
centreline of the test set (\PHPrunedMedShare\% of medium, \PHPrunedThickShare\% of thick).

\textit{Matched operating point.} For every run, the threshold on a grid of 0.01 to 0.99 (step 0.01) whose mean
validation precision was closest to a target precision was applied to the test set. The main target was the U-Net's mean
validation precision at its selected thresholds (\PHMatchTarget, mean over its three seeds); 0.90 and 0.95 were also used.

\textit{Resampling schemes.} Besides the image bootstrap of the main analysis, the test images were grouped into
possibly same-patient clusters by the rules used for the split (byte-identical image files, identical marking pattern with
consecutive numbers, isolated pairs with an identical marking pattern), applied to pairs within the test set, plus
identical masks: \PHNClusters{} clusters, \PHNMultiClusters{} of them with more than one image (\PHNClusteredImages{}
images, at most \PHLargestCluster{} per cluster). The cluster bootstrap resampled clusters. The seed-and-cluster bootstrap
also resampled, in each replicate, the three seeds of each model with replacement: jointly for comparisons among the four
CNNs, whose seeds are coupled, and independently for comparisons of SAM with the U-Net. Seed-paired values compute the
statistic from seed $s$ of every model ($s = 0, 1, 2$). All intervals are 95\% percentile intervals from \PHNBoot{}
replicates.

\textit{Official release.} The FIVES archive on figshare (version 1) was downloaded, its checksum compared with the one
listed by figshare, and its \OffNPng{} image and mask files compared with the copy used, by SHA-256 hash and, where the
files differed, by decoded pixels. All \OffBytesEqualImg{} image files were byte-identical. The masks differed only in
storage: the release stores three identical colour channels (\OffPixEqualMaskRGBA{} masks with an additional fully opaque
alpha channel) and the copy one channel, and all \OffPixEqualMask{} were pixel-identical after reduction to one channel.
The audit findings derived again from the official files were the same: \OffNDupGroups{} pairs of byte-identical images
with different masks, \OffNEmptyMasks{} empty masks, and Dice \OffSameEyeDice{} between the masks of the same-eye pair.
Images were only hashed and compared as arrays.

\subsection*{A10. Attention gates and data extraction}
Each additive attention gate~\cite{oktay2018attention} multiplies a skip connection $x$ by
\begin{equation}
\alpha = \sigma\left(\psi\left(\mathrm{ReLU}\left(\theta_x(x) + \phi_g(g)\right)\right)\right),
\end{equation}
where $g$ are the decoder features one level coarser, $\theta_x$ is a 2 $\times$ 2 convolution with stride 2 and no bias,
$\phi_g$ and $\psi$ are 1 $\times$ 1 convolutions, and the intermediate width is half the channels of $x$; $\alpha$ is
enlarged bilinearly and $\alpha x$ replaces $x$. When the images and masks were extracted from the Hugging Face copy, all
\NImages{} image and mask files written to disk were byte-identical to the bytes stored in the copy.
